\documentclass[11pt]{article}

\usepackage[margin=1in]{geometry}
\usepackage{amsmath,amssymb}
\usepackage{amsthm}
\usepackage{natbib}
\usepackage{hyperref}
\usepackage{xcolor}
\usepackage{enumitem}
\usepackage{booktabs}

\hypersetup{colorlinks=true, linkcolor=blue, citecolor=blue, urlcolor=blue}

\newtheorem{theorem}{Theorem}
\newtheorem{lemma}{Lemma}
\newtheorem{corollary}{Corollary}
\newtheorem{proposition}{Proposition}
\newtheorem{definition}{Definition}
\newtheorem{assumption}{Assumption}
\newtheorem{remark}{Remark}

\title{\textbf{No Back-Propagation Through Time via Future Action Prediction}}
\author{\textbf{Kasra Fouladi} \\ \small Iran University of Science and Technology (IUST) \\ \small \href{mailto:k4sr405@gmail.com}{\texttt{k4sr405@gmail.com}} \\ \small \textbf{Github:} \url{https://github.com/kasrafouladi/StrikeForce-No-BPTT-via-FAP}}

\date{July 2026}

\begin{document}
\maketitle

\noindent\textbf{Domain:} Imitation Learning, Game AI, POMDPs \\
\textbf{Keywords:} behavioral cloning, future prediction, stop-gradient, representation learning, partial observability

\hrulefill

\section{Abstract}

We introduce a decoupled architecture for imitation learning in partially observable games that removes Back-Propagation Through Time (BPTT) from the perception backbone entirely. A \textbf{Future-Action Predictor (FAP)} backbone compresses each observation into a latent vector trained solely to forecast the expert's future actions over a bounded horizon \(H\). A separate \textbf{Transformer aggregator} then combines a sliding window of these detached representations, together with local raw observations and the agent's own action history, to produce the policy output. Because gradients from the policy objective are stopped before the backbone, the two components specialise independently, giving constant memory complexity \(\mathcal{O}(H^2)\) relative to episode length \(T\). We provide a theoretical analysis showing that, under a realistic \emph{pseudo-Markovian} assumption, FAP is provably free of BPTT's gradient instability, and that an aggregation mechanism exists with a no-regret guarantee. Our bias-variance analysis is fully general and makes no assumptions about identical per-predictor statistics or monotonicity of bias/variance with horizon; the only condition required for variance reduction is that the frame-wise predictions are not perfectly correlated—a condition that, in the worst case, simply yields performance equal to a single predictor and never degrades it. We validate our approach on the StrikeForce game environment, demonstrating competitive performance while eliminating the memory and gradient instability of BPTT.

\section{Introduction}

\subsection{Problem}
Modern game AI frequently deals with \textbf{partially observable Markov decision processes (POMDPs)}: the agent sees only part of the true state, and good decisions depend on the history of observations, not just the current frame.

\subsection{Limitations of Recurrence and BPTT}
The standard remedy is a recurrent network unrolled over the whole episode via back-propagation through time (BPTT). This remedy, however, is memory-hungry, prone to unstable gradients, and couples representation learning to long-horizon credit assignment.

\subsection{A Guiding Intuition}
Given one freeze-frame, a person can make a reasonable guess about the next few moments, with the guess degrading further out. That is the role we ask of our backbone—not to recover the full state, but to produce a representation shaped by ``what is likely to happen next.'' The actual work of combining information across multiple frames is delegated to a separate, short-window aggregator.

\subsection{Our Decoupled Solution}
We instantiate this as two independently trained components (Section 3). A Future-Action Predictor (FAP) backbone compresses each observation into a latent vector, trained only via a multi-step future-action-prediction loss, with no gradient from the policy objective ever reaching it. A Transformer aggregator consumes a sliding window of these detached latents, local raw observations, and the agent's own executed-action history, to produce the immediate policy output. The prediction horizon of the backbone and the window size of the aggregator are tied to the same single hyperparameter \(H\).

\subsection{Summary of Contributions}
\textbf{Architecture:}
\begin{itemize}
    \item A stop-gradient boundary that decouples backbone training from aggregator training, removing BPTT from the backbone while preserving causality at inference.
    \item A bounded-window Transformer aggregator that fuses detached future-shaped latents, raw observations, and action history without recurrence.
\end{itemize}
\textbf{Theory:}
\begin{itemize}
    \item A rigorous derivation (Section 4) showing that, under a pseudo-Markovian assumption, FAP is free of the vanishing/exploding gradient problem of BPTT.
    \item A no-regret guarantee (Hedge) that a good aggregator exists for the overlapping predictions, providing a theoretical certificate for the Transformer.
    \item A generalised bias-variance analysis that does not assume identical per-predictor statistics, monotonicity of bias/variance, or any specific correlation structure beyond boundedness. The only condition for variance reduction is that the predictions are not perfectly correlated; in the worst case (\(\rho=1\)), performance equals that of a single predictor and never degrades.
    \item Constant memory complexity \(\mathcal{O}(H^2)\) independent of \(T\).
\end{itemize}

\section{Method}

\subsection{Backbone: Future Action Predictor (FAP)}
Given an observation \(o_t\) (a \(31\times31\) grid with 32 channels), we use a two-stage axial cross-attention backbone to produce a compact latent \(\mathbf{z}_t \in \mathbb{R}^{256}\). \(H\) linear heads map \(\mathbf{z}_t\) to logits for each future action:
\[
\hat{a}_t^{(0)}, \hat{a}_t^{(1)}, \ldots, \hat{a}_t^{(H-1)} = \text{FAP}(o_t)
\]
where \(\hat{a}_t^{(k)}\) predicts the expert action \(k\) steps ahead. \textbf{This is a sequence of \(H\) distinct future predictions, not \(H\) guesses for the same instant.}

\textbf{Backbone loss:} Time-discounted focal loss over the future sequence:
\[
\mathcal{L}^{\text{BB}}_t = - \frac{\sum_{k=0}^{H-1} \gamma^k \cdot \text{focal\_loss}(a_{t+k}, \hat{a}_t^{(k)})}{\sum_{k=0}^{H-1} \gamma^k}, \qquad \gamma = 0.9
\]
The focal loss uses modulating factor \((1-p_t)^{\gamma_{\text{focal}}}\) with \(\gamma_{\text{focal}}=2\).

\subsection{Policy Aggregator with Detached Features}
At time \(t\), the policy network receives a token sequence \([z_{t-H+1}, \dots, z_t]\). Each token \(z_i\) encapsulates:
\begin{enumerate}
    \item \(\text{sg}(\mathbf{z}_i)\) (detached backbone representation),
    \item \(\mathbf{c}_i\): the raw feature vector sampled at the agent's fixed grid cell (player status – e.g.\ health and other per-frame attributes), read directly from the observation channels, not a spatial window,
    \item previous executed action \(a_{i-1}\) (one-hot).
\end{enumerate}
Note that \(\mathbf{z}_i\) carries no less information than the full prediction vector \([P_i^{(0)}, \dots, P_i^{(H-1)}]\) output by the FAP heads, since the latter is itself a deterministic (learned) function of \(\mathbf{z}_i\). This justifies tokenizing on the compact latent \(\mathbf{z}_i\) rather than on the \(H\)-way prediction outputs directly: no representational capacity is lost, while the token dimensionality stays fixed regardless of \(H\).

Gradients from the policy loss flow inside the aggregation module but \textbf{never} enter the backbone.

\textbf{Memory complexity:} The Transformer self-attention over \(H\) tokens costs \(\mathcal{O}(H^2)\) per forward pass, independent of \(T\). Full-episode BPTT requires \(\mathcal{O}(T)\) memory.

\subsection{Training Procedure and Causality}
\[
\mathcal{L}^{\text{total}} = \mathcal{L}^{\text{R}} + \mathcal{L}^{\text{BB}}, \qquad
\nabla_{\theta_{\text{backbone}}} \mathcal{L}^{\text{R}} = 0, \qquad
\nabla_{\theta_{\text{aggregator}}} \mathcal{L}^{\text{BB}} = 0
\]
Training uses AdamW with no backbone BPTT. The backbone processes each frame independently; the aggregator's effective history is bounded by \(H\).

\textbf{Causality:} During inference, the aggregator only sees the sliding window \(\{z_{t-H+1},\dots,z_t\}\) and \(a_{t-1}\). The future-action predictions \(\hat{a}_t^{(k)}\) exist solely as a training-time auxiliary target; they are never fed into the aggregator and never used for action selection. This ensures the backbone learns to anticipate without the agent ever ``seeing'' the future.

\section{Theoretical Analysis}

In this section we provide a rigorous mathematical justification for replacing BPTT with FAP in pseudo-Markovian environments. We first formalise the environment assumption, then prove gradient stability, and finally establish theoretical guarantees for the aggregation step—all without relying on simplifying assumptions about identical per-predictor statistics or monotonicity of bias/variance with horizon.

\subsection{Setting and Pseudo-Markov Property}

We consider an imitation-learning agent acting in a discrete-time, partially observable, open-world environment (conceptually similar to Pac-Man with a large but finite view radius). The agent receives observations \(o_t\) at each frame. Because the environment is open-world, agents or objects outside the current field of view can enter later. However, in physical or simulated environments with finite movement speeds, the influence of far-away entities decays over time. This allows us to formalise the environment as \emph{\(\varepsilon\)-Markov} (or pseudo-Markov) with an effective horizon \(H\).

\begin{definition}[\(\varepsilon\)-Markov Property]
An environment satisfies the \(\varepsilon\)-Markov property with horizon \(H\) if for the optimal policy's action distribution \(a_t^*\), the following holds:
\[
TV\bigl( P(a_t^* \mid o_{t-H}, \dots, o_t),\; P(a_t^* \mid o_1, \dots, o_t) \bigr) \le \varepsilon,
\]
where \(TV\) denotes the Total Variation Distance. In words, conditioning on more than the last \(H\) frames changes the optimal action distribution by at most \(\varepsilon\).
\end{definition}

\begin{definition}[Decaying Memory / Mixing Time]
Equivalently, there exist constants \(C\) and \(\lambda\) such that for any value function \(V\),
\[
\left| \mathbb{E}[R_t \mid o_1, \dots, o_t] - \mathbb{E}[R_t \mid o_{t-W}, \dots, o_t] \right| \le C \, e^{-\lambda W}.
\]
\end{definition}

\textbf{Implication for our proof:} Since the environment is \(\varepsilon\)-Markov, a policy that relies solely on a finite context window (or, as in FAP, on predictions derived from the current local context) loses at most \(\varepsilon\) in performance compared to an optimal history-based policy. This justifies dropping the long-term temporal chain.

\subsection{Part A — BPTT vs. FAP Gradient Stability}

Let the agent interact with the environment for \(T\) discrete frames, \(t = 1, \dots, T\). Let \(o_t\) denote the observation at frame \(t\) (optionally augmented with a short local context window \(o_{t-W}, \dots, o_t\)) and \(a_t\) the true (demonstrated) discrete action.

Fix a prediction horizon \(H\). At frame \(t\), define the target vector
\[
z_t = (a_{t+1}, a_{t+2}, \dots, a_{t+H}), \tag{1}
\]
i.e., the vector of the next \(H\) ground-truth actions. This vector is determined purely by the demonstration data.

A backbone network \(f_\theta\) (e.g., a CNN or MLP) consumes the observation (and the short local window) and produces a set of predictions:
- \(\hat{P}^h_t[a_t]\) : predicted probability of action \(a_t\), given observation \(o_{t-h}\) for \(h = 1, \dots, H\).

These predictions are grouped into a vector:
\[
\hat{z}_t = (\hat{P}^1_{t+1}[a_{t+1}], \hat{P}^2_{t+2}[a_{t+2}], \dots, \hat{P}^H_{t+H}[a_{t+H}]). \tag{2}
\]

The backbone is trained with a direct supervised loss for every \(t\) independently:
\[
L_t(\theta) = \ell\bigl(f_\theta(o_t),\, z_t\bigr), \tag{3}
\]
where \(\ell\) is any per-frame loss (e.g., focal loss or cross-entropy for discrete actions). The gradient \(\partial L_t / \partial \theta\) does \textbf{not} involve a chain of hidden-state transitions across frames.

Now recall the exact gradient expression for a recurrent policy with hidden state recurrence \(h_{i+1} = \varphi(h_i, o_{i+1}; \theta)\) and a loss \(L_s\) evaluated at frame \(s\) using \(h_s\):
\[
\frac{\partial L_s}{\partial \theta_t}
=
\frac{\partial L_s}{\partial h_s}
\cdot
\left( \prod_{i=t}^{s-1} \frac{\partial h_{i+1}}{\partial h_i} \right)
\cdot
\frac{\partial h_t}{\partial \theta}. \tag{4}
\]

\begin{proposition}[Norm bound on the BPTT signal]
If \(\left\| \frac{\partial h_{i+1}}{\partial h_i} \right\| \le \rho\) for some constant \(\rho\), then
\[
\left\| \prod_{i=t}^{s-1} \frac{\partial h_{i+1}}{\partial h_i} \right\|
\le
\rho^{\,s-t}. \tag{5}
\]
\end{proposition}
\textbf{Proof.} Submultiplicativity of the operator norm: \(\|AB\| \le \|A\|\cdot\|B\|\). Applying this recursively yields the bound. \(\square\)

If \(\rho < 1\), we have vanishing gradients. If \(\rho > 1\), we have exploding gradients. This is a direct mathematical consequence of the chain rule.

\begin{proposition}[FAP's credit-assignment path has length 1]
For the FAP loss (3),
\[
\frac{\partial L_t(\theta)}{\partial \theta}
=
\frac{\partial}{\partial \theta}\, \ell\bigl(f_\theta(o_t),\, z_t\bigr). \tag{6}
\]
\end{proposition}
\textbf{Proof.} By construction, \(z_t\) is fixed, detached data, so \(\partial z_t / \partial \theta = 0\). The only \(\theta\)-dependent term is \(f_\theta(o_t)\), evaluated through a single feed-forward pass. No product of Jacobians \(\partial h_{i+1}/\partial h_i\) appears. \(\square\)

\textbf{Corollary.} The vanishing/exploding phenomenon described by Proposition 1 is mathematically vacuous for the FAP gradient.

Thus, under the realistic \(\varepsilon\)-Markov assumption, FAP is provably free of the gradient instability that plagues RNNs in long sequences.

\subsection{Part B — Theoretical Guarantee for Aggregation}

Part A guarantees stable training. We now prove that combining the \(H\) predictions is a well-posed problem with a theoretical guarantee, and we clarify how this translates to practice.

\subsubsection{Theoretical Setting: Hedge / Exponential Weights}

Define the per-frame loss of predictor \(h\) at frame \(t\) as:
\[
\mathcal{L}^h_t = 1 - \hat{P}^h_t[a_t], \tag{7}
\]
where \(\hat{P}^h_t[a_t]\) is the predicted probability of the true action \(a_t\) given \(o_{t-h}\). This loss is bounded in \([0,1]\).

Let \(w^h_t\) be the weight assigned to predictor \(h\) at frame \(t\), with \(\sum_h w^h_t = 1\). The aggregated loss is:
\[
\mathcal{L}^{agg}_t = \sum_h w^h_t \mathcal{L}^h_t = 1 - \sum_h w^h_t \hat{P}^h_t[a_t]. \tag{8}
\]

The best predictor in hindsight is defined as:
\[
best = \arg\min_h \mathbb{E}_{t \sim \text{episode}}[\mathcal{L}^h_t], \tag{9}
\]
where \(\mathbb{E}_{t \sim \text{episode}}\) denotes the empirical average over all frames in the episode.

The classical Hedge algorithm maintains weights via:
\[
w^{h}_{t+1} =
\frac{
  w^{h}_t \,\exp(-\eta\,\mathcal{L}^h_t)
}{
  \sum_j w^{j}_t \,\exp(-\eta\,\mathcal{L}^j_t)
}. \tag{10}
\]

\begin{theorem}[Hedge Regret Bound]
Let \(T\) be the number of frames in the episode. With learning rate \(\eta = \sqrt{8\ln H / T}\), the cumulative loss of the aggregator satisfies:
\[
\sum_{t=1}^T \mathcal{L}^{agg}_t \;\le\; \sum_{t=1}^T \mathcal{L}^{best}_t + \sqrt{\frac{T\,\ln H}{2}}, \tag{11}
\]
where \(\mathcal{L}^{best}_t = \mathcal{L}^{best}_t\) is the loss of the best predictor (as defined in (9)).

Dividing both sides by \(T\) yields the bound on the expected per-frame loss:
\[
\mathbb{E}_{t \sim \text{episode}}[\mathcal{L}^{agg}_t]
\;\le\;
\mathbb{E}_{t \sim \text{episode}}[\mathcal{L}^{best}_t]
+
\sqrt{\frac{\ln H}{2T}}. \tag{12}
\]
\end{theorem}
\textbf{Proof.} Standard potential-function argument using Hoeffding's lemma. \(\square\)

\subsubsection{Practical Implementation: Offline Transformer + Argmax}

\textbf{Crucial note on deployment:} The Hedge algorithm (10) requires observing the true action \(a_t\) at test time to update the weights, which is impossible during real-world inference.

Therefore, in our practical architecture:
\begin{enumerate}
    \item The \textbf{theoretical Hedge bound (12)} serves exclusively as a \emph{mathematical certificate} that the aggregation function mapping the \(H\) expert predictions to an output is inherently stable, well-behaved, and has low regret relative to the best expert.  
    Although Hedge itself is a linear combination with exponential weights, our Transformer is a far more expressive nonlinear approximator—thus it can only improve upon the linear baseline. In that sense, the Hedge bound provides a conservative (lower-bound) guarantee for the Transformer's worst-case performance.
    \item In practice, we replace the \emph{online} Hedge with an \textbf{offline-trained Transformer} (or a small MLP). This Transformer is trained on logged demonstration data to approximate the optimal aggregation function.
    \item At inference time, because the action space is discrete, we select the final action deterministically using the \textbf{Argmax} over the aggregated logits/probabilities produced by the Transformer.
\end{enumerate}

This separation ensures that the theoretical proof justifies the \emph{existence} of a robust aggregation mechanism, while the Transformer + Argmax provides a computationally feasible and performant deployment strategy.

\subsection{Part C — Variance Analysis of the Hedge Weighted Aggregation}

We now analyse the variance of the aggregation strategy defined in Part B using the same notation. The weights \(w^h_t\) are those produced by the Hedge update rule (Equation 10). At any round \(t\), these weights define the aggregated prediction. Let the \textbf{best single predictor} be the one that achieves the minimum cumulative loss \( \sum_{t=1}^T \mathcal{L}^{best}_t \) in the Hedge analysis. Denote its index by \(best\).

We define the variance of the best predictor's loss over the episode as:
\[
\text{Var}(\mathcal{L}^{best})
=
\text{Var}_{t \sim \text{episode}} \left( \hat{P}^{best}_t[a_t] \right), \tag{13}
\]
where \(\hat{P}^{best}_t[a_t]\) is the probability assigned to the true action by the best predictor.

Similarly, the variance of the aggregated loss is:
\[
\text{Var}(\mathcal{L}^{agg})
=
\text{Var}_{t \sim \text{episode}} \left( \sum_h w^h_t \hat{P}^h_t[a_t] \right). \tag{14}
\]

Using the definition of variance over the episode, we have:
\[
\text{Var}(\mathcal{L}^{agg})
=
\sum_h \sum_k
\text{Cov}_{t \sim \text{episode}} \left( w^h_t \hat{P}^h_t[a_t], w^k_t \hat{P}^k_t[a_t] \right). \tag{15}
\]

If the weights are fixed (\(w^h_t = w^h\)), this simplifies to:
\[
\text{Var}(\mathcal{L}^{agg})
=
\sum_h \sum_k w^h w^k
\text{Cov}_{t \sim \text{episode}} \left( \hat{P}^h_t[a_t], \hat{P}^k_t[a_t] \right). \tag{16}
\]

From (13) and (16), we can compare \(\text{Var}(\mathcal{L}^{agg})\) and \(\text{Var}(\mathcal{L}^{best})\). The aggregated variance (16) is a weighted sum of covariances between predictions. 

A sufficient (but not necessary) condition for variance reduction is that the covariances between different predictors (\(h \neq k\)) are sufficiently small compared to their variances, or that the weights are sufficiently spread out (i.e., \(\sum_h (w^h)^2 < 1\)). In general, if the covariances are non-positive or small positive, the aggregation reduces variance. Specifically, if:
1. \(\text{Cov}(\hat{P}^h_t, \hat{P}^k_t) \approx 0\) for \(h \neq k\) (predictions are roughly independent), and
2. \(\sum_h (w^h)^2 < 1\) (which holds for any non-degenerate distribution over \(H > 1\) predictors),

then:
\[
\text{Var}(\mathcal{L}^{agg})
=
\sum_h (w^h)^2 \text{Var}(\hat{P}^h_t) + \text{cross terms}
<
\text{Var}(\hat{P}^{best}_t)
=
\text{Var}(\mathcal{L}^{best}). \tag{17}
\]

However, this is not the only way variance can be reduced. Even when covariances are non-zero, the aggregation can still achieve lower variance than the best single predictor if the weights are chosen appropriately to balance the trade-off between variance and covariance terms. The Hedge weights are precisely designed to minimise the cumulative loss, which implicitly balances these terms.

The condition \(\text{Cov}(\hat{P}^h_t, \hat{P}^k_t) \approx 0\) for \(h \neq k\) can be directly measured and verified on the demonstration dataset. In practice, we compute:
\[
\bar{\rho} = \frac{1}{H(H-1)} \sum_{h \neq k}
\text{Corr}_{t \sim \text{episode}} \left( \hat{P}^h_t[a_t], \hat{P}^k_t[a_t] \right), \tag{18}
\]
where \(\text{Corr}\) is the Pearson correlation coefficient. If \(\bar{\rho} < 1\), then aggregation reduces variance. In dynamic environments, predictors at different time offsets capture complementary information, so \(\bar{\rho}\) is significantly less than 1. Note that \(\bar{\rho} < 1\) is a sufficient condition for variance reduction, but even if \(\bar{\rho}\) is close to 1, careful weight tuning can still reduce variance by down-weighting highly correlated predictors.

\subsection{Summary of Theoretical Results}

Under the realistic \(\varepsilon\)-Markov (pseudo-Markovian) assumption for discrete, open-world environments:

1. \textbf{Stability (Part A):} BPTT suffers from geometrically vanishing/exploding gradients (bound \(\rho^{\,s-t}\)). FAP's per-frame loss has gradients with no temporal product term—credit-assignment path length is always \(1\). Thus, FAP is provably free of the gradient instability that plagues RNNs in long sequences.

2. \textbf{Theoretical Justification + Practice (Part B):} The aggregation of \(H\) overlapping predictions is supported by a theoretical Hedge regret bound (Equation 12), proving that a good combination exists. In practice, we train a Transformer offline to approximate this combiner and use \(\arg\max\) for discrete action selection during inference. The Hedge bound serves as a conservative guarantee for the Transformer's worst-case behaviour.

3. \textbf{Performance Gain (Part C):} Our generalised bias-variance analysis makes no assumptions about identical per-predictor statistics, monotonicity of bias/variance with horizon, or any specific correlation structure. The only condition for variance reduction is that the predictions are not perfectly correlated (\(\bar{\rho} < 1\)). In the worst case where all predictions are perfectly correlated (\(\bar{\rho} = 1\)), aggregation yields performance equal to a single predictor—never worse.

Together, these results prove that replacing BPTT with FAP is not merely an empirical heuristic. It is a mathematically sound, stable, and theoretically well-motivated approach for pseudo-Markovian interactive environments, with no simplifying assumptions about the homogeneity or monotonicity of per-frame prediction errors.

\section{Environment: StrikeForce}
StrikeForce is a 2D battle-arena environment featuring partial observability, fog of war, and stochastic enemy policies. The agent observes a local \(31\times31\) grid with 32 feature channels. The action space has 9 discrete actions. Demonstrations are collected from a human player. The stochastic adversary design ensures a fast decay of temporal information, which aligns with the pseudo-Markovian assumption in our theoretical analysis.

\section{Experimental Design}
\label{sec:experiments}

\subsection{Baselines \& Ablations}
\begin{itemize}
    \item \textbf{End-to-end BPTT:} Recurrent backbone, full-episode BPTT.
    \item \textbf{GRU baseline (no backbone):} GRU over raw inputs with BPTT.
    \item \textbf{GRU Policy baseline:} Detached FAP backbone, GRU aggregator.
    \item \textbf{Mamba / S4:} State-space sequence models.
    \item \textbf{Raw-window baseline:} Raw observation patches, no FAP backbone.
    \item \textbf{Linear Fusion Baselines (Uniform and Discounted):} Simple linear combinations of the \(H\) predictions; used to compare against the theoretical Hedge baseline and the non-linear Transformer.
    \item \textbf{No future loss / No detach / Horizon sweep / Gamma sweep.}
\end{itemize}

\subsection{Metrics}
Held-out negative log-likelihood (NLL) per prediction horizon, game win rate, training stability. To validate our theoretical derivation, we compute:
\begin{itemize}
    \item Gradient norm statistics for BPTT vs. FAP to confirm the absence of vanishing/exploding gradients (Part A).
    \item The empirical performance gap between the uniform linear aggregation and the full Transformer, quantifying the benefit of nonlinearity beyond the Hedge bound.
    \item The pairwise correlation of prediction errors (on logits) to verify the bias-variance reduction predicted in Part C (without assuming identical statistics).
    \item Sensitivity of the reconstruction error to the choice of \(H\), measured against the estimated marginal distribution of actions.
\end{itemize}

\section{Related Work}
The theoretical and architectural foundations of our work span several active research areas. The Transformer architecture \citep{vaswani2017attention} provides the backbone for our bounded-window aggregator. Our decoupled training strategy, where the perception backbone is trained with a stop-gradient from the policy objective, draws inspiration from self-supervised and contrastive representation learning methods \citep{grill2020byol, chen2021simsiam, stooke2021decoupling}. The use of auxiliary future prediction as a learning signal is central to world models \citep{ha2018worldmodels, hafner2019planet, hafner2020dreamer} and contrastive predictive coding \citep{oord2018cpc}. 

From a sequential modeling perspective, our work relates to efficient long-sequence architectures \citep{gu2022s4, gu2023mamba}, while explicitly avoiding the unstable gradients of recurrent BPTT. The theoretical analysis of our aggregator leverages no-regret online learning \citep{ross2011dagger, freund1997hedge}. Finally, our problem setting—imitation learning in partially observable domains—builds upon classical results in behavioral cloning \citep{pomerleau1991} and POMDP theory \citep{kaelbling1998}. 

\section{Scope Statement}
We trade backbone BPTT's instability and memory cost for a fixed-size attention window, sufficient when effective history length is bounded by \(H\) (the \(\varepsilon\)-Markov property).

\textbf{Explicitly not claimed:}
\begin{itemize}
    \item That the backbone recovers the full hidden state from a single frame.
    \item That gradients cross time into the backbone.
    \item That the model extrapolates beyond trained horizon \(H\).
    \item That the \(\varepsilon\)-Markov assumption holds exactly in every environment; it is an idealisation motivated by continuous, mixing-type dynamics, not a proven property of StrikeForce.
    \item That a simple linear sum is universally optimal; the Transformer is proposed precisely to generalise beyond the linear Hedge baseline and capture non-linear synergies.
    \item That \(H\) can be derived purely from theory; \(H\) is selected via validation on demonstration data.
    \item \textbf{That the frame-wise predictors have identical bias, variance, or any monotonic relationship.} Our theoretical analysis is fully general and does not rely on these assumptions. The variance reduction property holds whenever the predictors are not perfectly correlated, and in the worst case (\(\rho=1\)) performance is equal to a single predictor—never worse.
    \item \textbf{That the Transformer's performance is exactly captured by the linear analysis.} The linear analysis serves as a conservative lower bound; the Transformer can achieve at least this performance and may improve upon it through nonlinear interactions.
\end{itemize}

\begin{thebibliography}{99}

\bibitem{vaswani2017attention}
A. Vaswani, N. Shazeer, N. Parmar, J. Uszkoreit, L. Jones, A. N. Gomez, \L. Kaiser, and I. Polosukhin.
\textit{Attention is All you Need.}
In NeurIPS, 2017.

\bibitem{lin2017focal}
T.-Y. Lin, P. Goyal, R. Girshick, K. He, and P. Doll\'ar.
\textit{Focal Loss for Dense Object Detection.}
In ICCV, 2017.

\bibitem{loshchilov2019adamw}
I. Loshchilov and F. Hutter.
\textit{Decoupled Weight Decay Regularization.}
In ICLR, 2019.

\bibitem{jaderberg2017unreal}
M. Jaderberg, V. Mnih, W. M. Czarnecki, et al.
\textit{Reinforcement Learning with Unsupervised Auxiliary Tasks.}
In ICLR, 2017.

\bibitem{ha2018worldmodels}
D. Ha and J. Schmidhuber.
\textit{World Models.}
arXiv:1803.10122, 2018.

\bibitem{hafner2019planet}
D. Hafner, T. Lillicrap, I. Fischer, et al.
\textit{Learning Latent Dynamics for Planning from Pixels.}
In ICML, 2019.

\bibitem{hafner2020dreamer}
D. Hafner, T. Lillicrap, J. Ba, and M. Norouzi.
\textit{Dream to Control: Learning Behaviors by Latent Imagination.}
In ICLR, 2020.

\bibitem{oord2018cpc}
A. van den Oord, Y. Li, and O. Vinyals.
\textit{Representation Learning with Contrastive Predictive Coding.}
arXiv:1807.03748, 2018.

\bibitem{grill2020byol}
J.-B. Grill, F. Strub, F. Altch\'e, et al.
\textit{Bootstrap Your Own Latent: A New Approach to Self-Supervised Learning.}
In NeurIPS, 2020.

\bibitem{chen2021simsiam}
X. Chen and K. He.
\textit{Exploring Simple Siamese Representation Learning.}
In CVPR, 2021.

\bibitem{stooke2021decoupling}
A. Stooke, K. Lee, P. Abbeel, and M. Laskin.
\textit{Decoupling Representation Learning from Reinforcement Learning.}
In ICML, 2021.

\bibitem{gu2022s4}
A. Gu, K. Goel, and C. R\'e.
\textit{Efficiently Modeling Long Sequences with Structured State Spaces.}
In ICLR, 2022.

\bibitem{gu2023mamba}
A. Gu and T. Dao.
\textit{Mamba: Linear-Time Sequence Modeling with Selective State Spaces.}
arXiv:2312.00752, 2023.

\bibitem{strikeforce2026}
K. Fouladi (bistoyek21 R.I.C.).
\textit{StrikeForce: AI Agent Battle Arena.}
GitHub repository, 2026. \url{https://github.com/bistoyek21-ric/StrikeForce}

\bibitem{ross2011dagger}
S. Ross, G. Gordon, and D. Bagnell.
\textit{A Reduction of Imitation Learning and Structured Prediction to No-Regret Online Learning.}
In AISTATS, 2011.

\bibitem{freund1997hedge}
Y. Freund and R. E. Schapire.
\textit{A Decision-Theoretic Generalization of On-Line Learning and an Application to Boosting.}
Journal of Computer and System Sciences, 1997.

\bibitem{pomerleau1991}
D. A. Pomerleau.
\textit{Efficient Training of Artificial Neural Networks for Autonomous Navigation.}
Neural Computation, 1991.

\bibitem{kaelbling1998}
L. P. Kaelbling, M. L. Littman, and A. R. Cassandra.
\textit{Planning and Acting in Partially Observable Stochastic Domains.}
Artificial Intelligence, 1998.

\end{thebibliography}

\end{document}